\documentclass[10pt]{article}
\usepackage[utf8]{inputenc}
\usepackage[T1]{fontenc}
\usepackage{amsmath}
\usepackage{amsfonts}
\usepackage{amssymb}
\usepackage[version=4]{mhchem}
\usepackage{stmaryrd}
\usepackage{bbold}
\usepackage{hyperref}
\hypersetup{colorlinks=true, linkcolor=blue, filecolor=magenta, urlcolor=cyan,}
\urlstyle{same}

\title{Constructing more quadratic APN functions with the QAM method }

\author{Yuyin Yu¹ • Léo Perrin²}
\date{}


%New command to display footnote whose markers will always be hidden
\let\svthefootnote\thefootnote
\newcommand\blfootnotetext[1]{%
  \let\thefootnote\relax\footnote{#1}%
  \addtocounter{footnote}{-1}%
  \let\thefootnote\svthefootnote%
}

%Overriding the \footnotetext command to hide the marker if its value is `0`
\let\svfootnotetext\footnotetext
\renewcommand\footnotetext[2][?]{%
  \if\relax#1\relax%
    \ifnum\value{footnote}=0\blfootnotetext{#2}\else\svfootnotetext{#2}\fi%
  \else%
    \if?#1\ifnum\value{footnote}=0\blfootnotetext{#2}\else\svfootnotetext{#2}\fi%
    \else\svfootnotetext[#1]{#2}\fi%
  \fi
}

\DeclareUnicodeCharacter{22C5}{\ifmmode\cdot\else{$\cdot$}\fi}

\begin{document}
\maketitle
Received: 19 October 2021 / Accepted: 17 June 2022 / Published online: 29 June 2022\\
© The Author(s), under exclusive licence to Springer Science+Business Media, LLC, part of Springer Nature 2022

\begin{abstract}
If used as S-boxes, APN functions provide optimal resilience against differential attacks. However, the very existence of APN permutations operating on an even number $n$ of bits (with $n \geq 8$ ) has been an open problem for nearly 30 years. A possible method to solve this problem consists in generating APN functions, and then exploring the CCZ-equivalence classes of these functions looking for a permutation. Following this goal, we found 5412 new quadratic APN functions on $\mathbb{F}_{2^{8}}$ using an approach based on so-called Quadratic APN Matrices (QAM). This brings the number of known CCZ-inequivalent APN functions on $\mathbb{F}_{2^{8}}$ to 26525 . Unfortunately, none of these new functions are CCZ-equivalent to permutations. A complete list (to the best of our knowledge) of known quadratic APN functions, including our new ones, has been added to sboxu for ease of study by others. In this paper, we recall how to construct new QAMs from a known one. Based on these results and on others on smaller fields, we make two conjectures: that the total number of CCZinequivalent quadratic APN functions on $\mathbb{F}_{2^{8}}$ exceeds 50000 , and that the full list of quadratic APN functions could be obtained by modifying only a small number of entries of the QAM, though such a search remains computationally infeasible at this stage. Finally, we propose a new model which can handle the last two columns together and avoid some redundant computation.
\end{abstract}

Keywords APN ⋅ QAM ⋅ ortho-derivative

Mathematics Subject Classification (2010) 06E30 $\cdot$ 11T06 $\cdot$ 94A60

\footnotetext{This article belongs to the Topical Collection: Boolean Functions and Their Applications VI Guest Editors: Lilya Budaghyan, Claude Carlet, Tor Helleseth, and Cunsheng Ding

Yuyin Yu\\
\href{mailto:yuyuyin@163.com}{yuyuyin@163.com}\\
Léo Perrin\\
\href{mailto:leo.perrin@inria.fr}{leo.perrin@inria.fr}\\
${ }^{1}$ School of Mathematics and Information Sciences, Guangzhou University, 510006 Guangzhou, China\\
${ }^{2}$ Inria, Paris, France
}\section*{1 Introduction}
Let $\mathbb{F}_{2^{n}}$ be the finite field with $2^{n}$ elements, for some integer $n$. A function $F: \mathbb{F}_{2^{n}} \rightarrow \mathbb{F}_{2^{n}}$ is said to be $A P N$ if and only if the following equation has at most 2 solutions $x$ for all $a \neq 0$ and for all $b$ in $\mathbb{F}_{2^{n}}$ :

$$
F(x+a)+F(x)=b
$$

Finding APN functions is easy. For example, the monomial $x \mapsto x^{3}$ is always APN, and it is also an APN permutation for odd values of $n$. However, finding APN permutations for $n$ even is less trivial. In fact, when $n$ is even, this task is an open problem which has been nicknamed the Big APN Problem. For $n=4$, we know that there is no such permutation since it is possible to enumerate all possible 4-bit permutations (up to a relevant form of equivalence). While we know an example of a 6-bit APN permutation (which we discuss below), whether such permutations exist for $n \geq 8, n$ even, is still an open problem.

Browning et al. [5] found the first APN permutation in dimension 6. Their idea was to check the CCZ-equivalence [8] class of a quadratic APN function: indeed, if an APN function is CCZ-equivalent to a permutation, then that permutation has to be APN. Browning et al. then provided a method to find APN permutations. Their idea was to check whether there are permutations in the CCZ-equivalence classes of APN functions. In this paper, we focus on how to construct quadratic APN functions in dimension 8, since 8 bit is the most popular structure in designing encryption algorithms. Edel and Pott [9] listed 23 CCZ-inequivalent APN functions on $\mathbb{F}_{2^{8}}$. Weng et al. [12] and Yu et al. [14] extended the length of the list to 8190 . Taniguchi [11] found 2 new ones and the total number was extended to 8192 . A very recent breakthrough was achieved by Beierle and Leander [1, 3], where 12921 new quadratic APN functions were found in dimension 8 . In total, 21113 CCZ-inequivalent quadratic APN functions in dimension 8 were known before this paper. We present another 5412 new quadratic APN functions. Thus, the number of CCZ-inequivalent quadratic APN functions in dimension 8 increases to 26525 (up to April 30, 2021). Readers can turn to [2, 7] for a comprehensive discussion of APN related topics,

We will recall how to modify a QAM to get some new QAMs in the sequel. The related theory and algorithm can be found in [14]. A discussion of our results on 8 bits, including conjectures about 8 -bit APN functions, are presented in Section 3. In Section 4 we will show why the QAM method is still valuable even now. Finally we propose a new model which can avoid some redundant computation in Section 5.

\section*{2 Notation}
The following notations and results are needed to understand our work.

\begin{itemize}
  \item Let $\eta_{1}, \eta_{2}, \ldots, \eta_{m}$ be $m$ elements on $\mathbb{F}_{2^{n}}(m, n \geq 1)$, and $B=\left(\eta_{1}, \eta_{2}, \ldots, \eta_{m}\right) \in \mathbb{F}_{2^{n}}^{m}$. Then $\operatorname{Span}(B)=\operatorname{Span}\left(\eta_{1}, \eta_{2}, \ldots, \eta_{m}\right)$ denotes the subspace spanned by $\left\{\eta_{1}, \eta_{2}, \ldots, \eta_{m}\right\}$ over $\mathbb{F}_{2}$. Further, $\operatorname{Rank}_{\mathbb{F}_{2}}(B)=\operatorname{Rank}_{\mathbb{F}_{2}}\left\{\eta_{1}, \eta_{2}, \ldots, \eta_{m}\right\}$ denotes the dimension of $\operatorname{Span}(B)$.
  \item Let $F(x)=\sum_{1 \leq t<i \leq n} c_{i, t} x^{i^{i-1}+2^{t-1}} \in \mathbb{F}_{2^{n}}[x]$ be a quadratic homogeneous polynomial (quadratic polynomial without linear and constant terms), then the coefficient matrix\\
$C_{F}$ is an $n \times n$ matrix such that $C_{F}[t, i]=C_{F}[i, t]=c_{i, t}$ for $1 \leq t<i \leq n$ and $C_{F}[i, i]=0$ for $1 \leq i \leq n$.
  \item $\delta(F)=\max _{a \in \mathbb{F}_{2^{n}}^{\star}, b \in \mathbb{F}_{2^{n}}}\left|\left\{x \in \mathbb{F}_{2^{n}}: F(x+a)+F(x)=b\right\}\right|$ denotes the differential uniformity of $F$.
\end{itemize}

Suppose that $\alpha=\left\{\alpha_{1}, \alpha_{2}, \ldots, \alpha_{n}\right\}$ is a basis of $\mathbb{F}_{2^{n}}$ over $\mathbb{F}_{2}$, and let $M_{\alpha} \in \mathbb{F}_{2^{n}}^{n \times n}$ with $M_{\alpha}[i, u]=\alpha_{u}^{2^{i-1}}$ for $1 \leq u, i \leq n$. The transpose of $M_{\alpha}$ is denoted by $M_{\alpha}^{t}$. For any quadratic homogeneous polynomial $F(x)$, if $H=M_{\alpha}^{t} C_{F} M_{\alpha}$, then $H$ is a symmetric matrix over $\mathbb{F}_{2^{n}}$ with zero main diagonal. Note that $\alpha$ can be any basis, while in our algorithm, we choose the normal basis to construct the matrix $M_{\alpha}$ for simplicity. Suppose

$$
\alpha=\left\{\alpha_{1}, \alpha_{2}, \ldots, \alpha_{n}\right\}=\left\{\gamma, \gamma^{2}, \ldots, \gamma^{2^{n-1}}\right\}
$$

is a normal basis of $\mathbb{F}_{2^{n}}$. Then we have $M_{\alpha}[i, u]=\gamma^{2^{i+u-2}}$ for $1 \leq u, i \leq n$. Specifically, we let $\gamma=g^{11}$ on $\mathbb{F}_{2^{8}}$, where $g$ is the default primitive element ( $g^{8}+g^{4}+g^{3}+g^{2}+1=0$ ) in Magma [4].

Our method for generating QAMs relies on the following concept from [14].

Definition 1 (14 QAM) Let $H=\left(h_{u, v}\right)_{n \times n}$ be an $n \times n$ matrix defined on $\mathbb{F}_{2^{n}}$. The matrix $H$ is called a Quadratic APN Matrix (QAM) if:

\begin{enumerate}
  \item $H$ is symmetric and the elements in its main diagonal are all zeros, and
  \item every nonzero linear combination of the $n$ rows of $H$ has rank $n-1$ over $\mathbb{F}_{2}$.
\end{enumerate}

Crucially, there is a one-to-one correspondence between quadratic homogeneous APN functions and a subset of such matrices, as explained by the following theorem from the same paper.

Theorem 1 (Theorem 1 of 14) Let $F(x)=\sum_{1 \leq t<i \leq n} c_{i, t} x^{2^{i-1}+2^{t-1}} \in \mathbb{F}_{2^{n}}[x], C_{F}$ and $M_{\alpha}$ be defined as above. Let


\begin{equation*}
H=M_{\alpha}^{t} C_{F} M_{\alpha} \tag{1}
\end{equation*}


Then, $\delta(F)=2^{k}$ if and only if any nonzero linear combination of the $n$ rows of $H$ has rank at least $n-k$ over $\mathbb{F}_{2}$. In particular, $F(x)$ is an APN polynomial on $\mathbb{F}_{2^{n}}$ if and only if $H$ is a QAM. In fact, (1) describes a one to one correspondence between quadratic homogeneous APN functions and QAMs.

\section*{3 New 8-bit quadratic APN functions}
\subsection*{3.1 Our results}
Using the search algorithm from [14], we could obtain 6794 APN functions by modifying a very small part (less than $0.5 \%$ ) of the last two columns of the corresponding QAM of $x^{3}$, and it takes about 3 years with a personal computer to get these functions. However, these do not all correspond to new CCZ-classes. In order to partition this set into CCZ-classes,\\[0pt]
we used a classical method based on CCZ-class invariants. Those we considered are listed in Section 3.2, the main one being based on the ortho-derivative [6].

In total, we have obtained 5412 new CCZ-classes of quadratic APN functions operating on 8 bits. These functions have been added to sboxu. ${ }^{1}$ The function

\begin{verbatim}
sboxU.known_functions.eightBitAPN.second_QAMs()
\end{verbatim}

returns a list containing their look-up tables. The function

\begin{verbatim}
sboxU.known_functions.eightBitAPN.all_quadratics()
\end{verbatim}

now also returns them along with all other known quadratic APN functions in this dimension, i.e. a total of 26525 functions.

\subsection*{3.2 Using class invariants}
In order to sort our functions into distinct EA-equivalence classes (and thus distinct CCZclasses), ${ }^{2}$ we used the approach based on the ortho-derivative introduced in [6]. First let us recall the definition of the ortho-derivative.

Definition 2 Let $F: \mathbb{F}_{2^{n}} \rightarrow \mathbb{F}_{2^{n}}$ be a quadratic APN function, and let $x \cdot y$ denote a scalar product of $x$ and $y$ (where $x$ and $y$ are in $\mathbb{F}_{2^{n}}$ ). Then the ortho-derivative of $F$ is the unique function $\pi_{F}: \mathbb{F}_{2^{n}} \rightarrow \mathbb{F}_{2^{n}}$ such that $\pi_{F}(0)=0, \pi_{F}(a) \neq 0$ if $a \neq 0$, and such that

$$
\pi_{F}(a) \cdot(F(x+a)+F(x)+F(a)+F(0))=0
$$

for all $a \in \mathbb{F}_{2^{n}}^{*}$ and all $x \in \mathbb{F}_{2^{n}}$.\\[0pt]
The crucial fact behind the sorting approach presented in [6] is that if two quadratic APN functions are EA-equivalent, then their ortho-derivatives are affine-equivalent. As a consequence, they need to have identical differential and extended Walsh spectra. Thus, two functions for which these spectra do not match cannot be EA-equivalent, and as a consequence cannot be CCZ-equivalent [13]. For each function $F$, we computed its orthoderivative $\pi_{F}$. Then, we sorted all functions $F$ according to the extended Walsh and differential spectra of their ortho-derivatives, keeping only one function when several had the same spectra. Here are some observations about these 6794 functions that we found.

\begin{itemize}
  \item There is a pair of spectra shared by 3 different functions, 245 pairs that are each shared by 2 different functions, and 6301 pairs that each have only 1 function. As a consequence, we can only prove that there are at least $6547(1+245+6301)$ distinct CCZclasses in the set we generated.
  \item Among these 6547 functions, only 2 had a pair of spectra that was already present in the set identified by Beierle and Leander [3]; and 1133 had already been found using the QAM method [14]. The intersections are distinct: there was no function among the ones that we found that was in either the data set of Beierle and Leander or the previous QAM one.
\end{itemize}

In the end, we can conclude that we found 5412 new EA-equivalence classes of quadratic APN functions.

\footnotetext{${ }^{1}$ \href{https://github.com/lpp-crypto/sboxU}{https://github.com/lpp-crypto/sboxU} (Updated in 2021-04-27, see the file "\href{http://secondQAM.py}{secondQAM.py}").\\
${ }^{2}$ Two quadratic APN functions are CCZ-equivalent if and only if they are EA-equivalent [13].
}Unfortunately, none of our new functions yield a new extended Walsh spectrum. More precisely, for a fixed function $F$, if we let $N_{i}$ denote the number of pairs $(a, b) \in \mathbb{F}_{2^{8}} \times \mathbb{F}_{2^{8}}$ with $b \neq 0$ such that $\left|\sum_{x \in \mathbb{F}_{2^{8}}}(-1)^{a \cdot x+b \cdot F(x)}\right|=i$, then we observe the following spectra:

$$
\begin{aligned}
& \mathcal{W}_{1}=\left\{N_{0}=16320, N_{16}=43520, N_{32}=5440\right\} \\
& \mathcal{W}_{2}=\left\{N_{0}=15600, N_{16}=44544, N_{32}=5120, N_{64}=16\right\} \\
& \mathcal{W}_{3}=\left\{N_{0}=14880, N_{16}=45568, N_{32}=4800, N_{64}=32\right\},
\end{aligned}
$$

where $\mathcal{W}_{1}$ occurs 5084 times, $\mathcal{W}_{2} 324$ times, and $\mathcal{W}_{3}$ only 4 times.\\
The $\Delta$ - and $\Gamma$-ranks are well known CCZ-class invariants defined as follows. First, let $S$ be any subset of $\mathbb{F}_{2^{n}} \times \mathbb{F}_{2^{n}}$. We denote by $r(S)$ the rank of the binary matrix $M(S)$ of dimension $2^{n} \times 2^{n}$ which is such that $M_{x, y}=1$ if and only if $x+y \in S$. Then, we have that the $\Gamma$-rank of a function $F: \mathbb{F}_{2^{n}} \rightarrow \mathbb{F}_{2^{n}}$ is equal to

$$
\Gamma_{F}=r\left(\left\{(x, F(x)): x \in \mathbb{F}_{2^{n}}\right\}\right),
$$

and that its $\Delta$-rank is

$$
\Delta_{F}=r\left(\left\{(a, b) \in \mathbb{F}_{2^{8}} \times \mathbb{F}_{2^{8}}: F(x+a)+F(x)=b \text { has at least } 1 \text { solution }\right\}\right) .
$$

However, the $\Delta$-rank and $\Gamma$-rank are not distinguishable enough, since the 6794 APN functions can only be divided into 9 classes.\\
$(452,14044),(452,14046),(452,14048),(452,14050),(454,14042),(454,14044)$, $(454,14046),(454,14048),(454,14050)$.

The number of APN functions in each class is\\
$7,21,45,6,31,482,2363,2831,1008$.\\
It takes about 1500 hours to compute their $\Delta$-ranks and $\Gamma$-ranks, ${ }^{3}$ while it takes 21 seconds to compute the extended Walsh and differential spectra of their ortho-derivatives.

We have also computed the multisets $\Sigma_{4}^{F}(0)$ for each $F$ in our set of functions, where this quantity is the multiset defined as follows [10]:

$$
\Sigma_{4}^{F}(0)=\left\{\sum_{i=0}^{3} F\left(x_{i}\right):\left(x_{0}, \ldots, x_{3}\right) \in\left(\mathbb{F}_{2^{n}}\right)^{4}, \text { and } \sum_{i=0}^{3} x_{i}=0\right\}
$$

and is an EA-class invariant. We have implemented an optimized version of this invariant, and added the corresponding function to sboxu. It is called sigma\_multiplicities, and relies on the Walsh coefficients of $F$ to compute $\Sigma_{4}^{F}(0)$.

Evaluating this invariant on our full set of function turned out a bit slower than the ortho-derivative (it takes several minutes), and 4655 distinct values were found. This invariant is thus far more useful than all but the ortho-derivative: it can be computed significantly faster, and provides much more differentiation. It also has the significant advantage over the ortho-derivative that it does not require the function investigated to be both quadratic and APN, these requirements limiting the usefullness of the ortho-derivative to work such as ours or the one of Beierle and Leander [3].

\subsection*{3.3 Some conjectures}
At this point, the total number of CCZ-inequivalent quadratic APN functions on $\mathbb{F}_{2^{8}}$ is higher than 26000. However, we believe that this number is still far from complete.

\footnotetext{${ }^{3}$ with codes from \href{https://boolean.h.uib.no/mediawiki/index.php/Magma_Code}{https://boolean.h.uib.no/mediawiki/index.php/Magma\_Code}.
}Conjecture 1 The total number of CCZ-inequivalent quadratic APN functions on $\mathbb{F}_{2^{8}}$ is higher than 50000.

We list some facts to support Conjecture 1 that correspond to the experiments we made by looping through all the QAMs with a given structure for a given $n$.\\
(1) In dimension 8, we can still construct a quadratic APN function every 24 hours with the QAM method, and there is an about $79 \%$ probability that it is new compared to all known ones.\\
(2) In dimension 7, when $230\left(47 \%=\frac{230}{488}\right.$ of the total number) CCZ-inequivalent quadratic APN functions have been found, there is an about $79 \%$ probability that the next APN function constructed by the QAM method is new (i.e. not among the first 230).\\
(3) In dimension 6, when $6\left(46 \%=\frac{6}{13}\right.$ of the total number) CCZ-inequivalent quadratic APN functions have been traversed, there is an about $75 \%$ probability that the next APN function constructed by the QAM method is not among the first 6 found.

Based on the above facts, we guess that the total number of CCZ-inequivalent quadratic APN functions in dimension 8 is at least twice the number of the known ones.

The following conjecture may provide a method to construct the complete list of quadratic APN functions in dimension 8.

Conjecture 2 Let $C_{F}=M_{\alpha}^{-1} H_{8}\left(M_{\alpha}^{t}\right)^{-1}$ be the coefficient matrix of $F(x) \in \mathbb{F}_{2^{8}}[x]$, where $H_{8}$ is such that

$$
H_{8}=\left(\begin{array}{cccccccc}
0 & g^{34} & g^{81} & g^{83} & g^{170} & g^{106} & \mathbf{x}_{13} & \mathbf{x}_{7} \\
g^{34} & 0 & g^{68} & g^{162} & g^{166} & g^{85} & \mathbf{x}_{12} & \mathbf{x}_{6} \\
g^{81} & g^{68} & 0 & g^{136} & g^{69} & g^{77} & \mathbf{x}_{11} & \mathbf{x}_{5} \\
g^{83} & g^{162} & g^{136} & 0 & g^{17} & g^{138} & \mathbf{x}_{10} & \mathbf{x}_{4} \\
g^{170} & g^{166} & g^{69} & g^{17} & 0 & g^{34} & \mathbf{x}_{9} & \mathbf{x}_{3} \\
g^{106} & g^{85} & g^{77} & g^{138} & g^{34} & 0 & \mathbf{x}_{8} & \mathbf{x}_{2} \\
\mathbf{x}_{13} & \mathbf{x}_{12} & \mathbf{x}_{11} & \mathbf{x}_{10} & \mathbf{x}_{9} & \mathbf{x}_{8} & 0 & \mathbf{x}_{1} \\
\mathbf{x}_{7} & \mathbf{x}_{\mathbf{6}} & \mathbf{x}_{\mathbf{5}} & \mathbf{x}_{4} & \mathbf{x}_{3} & \mathbf{x}_{2} & \mathbf{x}_{1} & 0
\end{array}\right) .
$$

All CCZ-inequivalent classes of quadratic APN functions on $\mathbb{F}_{2^{8}}$ can be obtained by letting $\mathrm{x}_{1}, \mathrm{x}_{2}, \ldots, \mathrm{x}_{12}$ and $\mathrm{x}_{13}$ traverse $\mathbb{F}_{2^{8}}$.

The left upper $6 \times 6$ submatrix of $\mathrm{H}_{8}$ is the same as the one of the corresponding QAM of $x^{3}$, and all known QAM-based APN functions are constructed by modifying the last two columns (and rows) of the matrix $\mathrm{H}_{8}$. However, we have only traversed less than $1 \%$ elements of the last two columns (and rows), and there is an about $79 \%$ probability that the next QAM-based APN function is new compared to all known ones.

We had to generate more than 200 (about $16 \times 13$ ) and more than 3000 (about $8 \times 488$ ) quadratic APN functions using the QAM method in order to obtain the complete list of quadratic APN functions in dimension 6 and dimension 7, respectively. Therefore, we may need to generate more than $200000(4 \times 50000)$ quadratic APN functions in order to get the complete list in dimension 8 using a QAM-based approach. Thus, there is a high probability that we can get the full list after traversing $x_{1}, x_{2}, \ldots, x_{12}$ and $x_{13}$, since using this method we can construct at least 2000000 quadratic APN functions. Unfortunately, this would require substantial computations corresponding to centuries of CPU time.

\section*{4 Avoiding EA-equivalent APN functions}
In this section, we give a detailed description of how to avoid traversing some EA-equivalent APN functions in Conjecture 2. We need the following theorem to prove Proposition 3.

Theorem 2 (Theorem 2 of 14) Let $H \in \mathbb{F}_{2^{n}}^{n \times n}$ be a symmetric matrix with main diagonal elements all zeros, and $P \in \mathbb{F}_{2}^{n \times n}$ be an invertible matrix. Suppose $H^{\prime}=P^{t} H P$, then the quadratic functions defined by $H$ and $H^{\prime}$ are EA-equivalent. Especially, $H$ is a $Q A M$ if and only if $H^{\prime}$ is a $Q A M$.

Proposition 3 Suppose that $H_{8}$ is the same as in Conjecture 2. Let

$$
V=\operatorname{Span}\left(g^{34}, g^{81}, g^{83}, g^{170}, g^{106}\right)
$$

Then we can partition $\mathbb{F}_{2^{8}}$ into

$$
V \cup\left(V+a_{1}\right) \cup\left(V+a_{2}\right) \cup\left(V+a_{3}\right) \cup\left(V+a_{4}\right) \cup\left(V+a_{5}\right) \cup\left(V+a_{6}\right) \cup\left(V+a_{7}\right)
$$

for some $a_{1}, a_{2}, a_{3}, a_{4}, a_{5}, a_{6}$ and $a_{7}$. If $H_{8}$ is a $Q A M$, then in the sense of EA-inequivalence, we only need to consider the set $A=\left\{a_{1}, a_{2}, a_{3}, a_{4}, a_{5}, a_{6}, a_{7}\right\}$ when choosing $x_{13}$. Similarly, in the sense of EA-inequivalence, the cardinalities of the candidate set for $x_{12}, x_{7}, x_{6}$ are $2^{7} -2^{4}, 3$ and $2^{7}-2^{5}$, respectively.

Proof Let $H_{8}$ be a QAM, then we have $\operatorname{Rank}_{\mathbb{F}_{2}}\left(g^{34}, g^{81}, g^{83}, g^{170}, g^{106}, x_{13}\right)=6$, which implies that $x_{13} \in \mathbb{F}_{2^{8}} \backslash V$. Suppose $x_{13}=\lambda_{1} g^{34}+\lambda_{2} g^{81}+\lambda_{3} g^{83}+\lambda_{4} g^{170}+\lambda_{5} g^{106}+a \left(\lambda_{i} \in \mathbb{F}_{2}\right.$ and $\left.a \in A\right)$, then we have

$$
H_{8}^{\prime}=P^{t} H_{8} P=\left(\begin{array}{cccccccc}
0 & g^{34} & g^{81} & g^{83} & g^{170} & g^{106} & \mathbf{a} & \mathbf{x}_{\mathbf{7}} \\
g^{34} & 0 & g^{68} & g^{162} & g^{166} & g^{85} & \mathbf{y}_{\mathbf{2}} & \mathbf{x}_{\mathbf{6}} \\
g^{81} & g^{68} & 0 & g^{136} & g^{69} & g^{77} & \mathbf{y}_{\mathbf{3}} & \mathbf{x}_{\mathbf{5}} \\
g^{83} & g^{162} & g^{136} & 0 & g^{17} & g^{138} & \mathbf{y}_{\mathbf{4}} & \mathbf{x}_{\mathbf{4}} \\
g^{170} & g^{166} & g^{69} & g^{17} & 0 & g^{34} & \mathbf{y}_{\mathbf{5}} & \mathbf{x}_{\mathbf{3}} \\
g^{106} & g^{85} & g^{77} & g^{138} & g^{34} & 0 & \mathbf{y}_{\mathbf{6}} & \mathbf{x}_{\mathbf{2}} \\
\mathbf{a} & \mathbf{y}_{\mathbf{2}} & \mathbf{y}_{\mathbf{3}} & \mathbf{y}_{\mathbf{4}} & \mathbf{y}_{\mathbf{5}} & \mathbf{y}_{\mathbf{6}} & 0 & \mathbf{x}_{\mathbf{1}} \\
\mathbf{x}_{7} & \mathbf{x}_{\mathbf{6}} & \mathbf{x}_{\mathbf{5}} & \mathbf{x}_{\mathbf{4}} & \mathbf{x}_{\mathbf{3}} & \mathbf{x}_{\mathbf{2}} & \mathbf{x}_{\mathbf{1}} & 0
\end{array}\right),
$$

where $y_{i}=\lambda_{1} H_{8}[i, 2]+\lambda_{2} H_{8}[i, 3]+\lambda_{3} H_{8}[i, 4]+\lambda_{4} H_{8}[i, 5]+\lambda_{5} H_{8}[i, 6]+x_{14-i}$ for any $2 \leq i \leq 6$, and

$$
P=\left(\begin{array}{llllllll}
1 & 0 & 0 & 0 & 0 & 0 & 0 & 0 \\
0 & 1 & 0 & 0 & 0 & 0 & \lambda_{1} & 0 \\
0 & 0 & 1 & 0 & 0 & 0 & \lambda_{2} & 0 \\
0 & 0 & 0 & 1 & 0 & 0 & \lambda_{3} & 0 \\
0 & 0 & 0 & 0 & 1 & 0 & \lambda_{4} & 0 \\
0 & 0 & 0 & 0 & 0 & 1 & \lambda_{5} & 0 \\
0 & 0 & 0 & 0 & 0 & 0 & 1 & 0 \\
0 & 0 & 0 & 0 & 0 & 0 & 0 & 1
\end{array}\right) .
$$

According to Theorem 2, if $H_{8}$ is a QAM, then $H_{8}^{\prime}$ is also a QAM, and the corresponding APN functions of $H_{8}$ and $H_{8}^{\prime}$ are EA-equivalent. Therefore, we only need to consider $x_{13} \in A$ when choosing it.

Similarly, let $U=\operatorname{Span}\left(g^{34}, g^{68}, g^{162}, g^{166}, g^{85}\right)$, and $B \cup\left(B+g^{34}\right)$ be a partition of $\mathbb{F}_{2^{8}} \backslash U$, and we only need to consider $x_{12} \in B$ when traversing it.

When the values of $x_{13}$ and $x_{12}$ have been chosen, then let $E=\operatorname{Span}\left(g^{34}, g^{81}, g^{83}, g^{170}, g^{106}\right.$, $\left.x_{13}\right), F=\operatorname{Span}\left(g^{34}, g^{68}, g^{162}, g^{166}, g^{85}, x_{12}\right), E \cup\left(E+c_{1}\right) \cup\left(E+c_{2}\right) \cup\left(E+c_{3}\right)$ be a partition of $\mathbb{F}_{2^{8}}$, and $G \cup\left(G+g^{34}\right)$ be a partition of $\mathbb{F}_{2^{8}} \backslash F$. Then we only need to consider $x_{7} \in\left\{c_{1}, c_{2}, c_{3}\right\}$ and $x_{6} \in G$ when choosing them.

Note that $|A|=7,|B|=2^{7}-2^{4},\left|\left\{c_{1}, c_{2}, c_{3}\right\}\right|=3$ and $|G|=2^{7}-2^{5}$.

Note that If $H$ is a QAM, then adding column $i$ and row $i$ to column $j$ and row $j$ respectively will produce another QAM $H^{\prime}$, and the corresponding functions of $H$ and $H^{\prime}$ are EA-equivalent. Proposition 3 provides a method to avoid such EA-equivalence in searching. The results in Proposition 3 can be generalized to any dimension. Based on it, we can significantly decrease the number of candidate values when traversing it. According to our observation, the cardinalities of the candidate values for $x_{i}(1 \leq i \leq 13)$ have only a small fluctuation when running Algorithm 1 of [14]. Thus, we can estimate the total number of the candidate values for all $x_{i}(1 \leq i \leq 13)$.

Conjecture 3 Suppose $\mathrm{H}_{8}$ is the same as in Conjecture 2. Then the complexity of the QAM-based method for traversing all $\mathrm{x}_{\mathrm{i}}(1 \leq \mathrm{i} \leq 13)$ is not higher than $2^{50}$.

According to our experimental results, the cardinalities of the candidate values for $x_{2}$ and $x_{1}$ are either 1 or 0 , so we neglect them in the following. Based on Proposition 3, the cardinalities of the candidate values for $x_{13}$ is 7 . For each $a \in A$, we choose one cardinality combination for

$$
x_{13} \times x_{12} \times x_{11} \times x_{10} \times x_{9} \times x_{8} \times x_{7} \times x_{6} \times x_{5} \times x_{4} \times x_{3}
$$

Then we can get either:

\begin{itemize}
  \item $7 \times 100 \times 154 \times 87 \times 20 \times 2 \times 3 \times 64 \times 90 \times 29 \times 3=563976483840000 \approx 2^{49.0}$,
  \item $7 \times 96 \times 147 \times 89 \times 26 \times 2 \times 3 \times 72 \times 78 \times 31 \times 4=955102533525504 \approx 2^{49.8}$,
  \item $7 \times 96 \times 147 \times 78 \times 32 \times 2 \times 3 \times 72 \times 84 \times 27 \times 4=966313344761856 \approx 2^{49.8}$,
  \item $7 \times 96 \times 152 \times 79 \times 29 \times 3 \times 3 \times 72 \times 92 \times 28 \times 2=781247805456384 \approx 2^{49.5}$,
  \item $7 \times 96 \times 154 \times 85 \times 23 \times 4 \times 3 \times 72 \times 92 \times 23 \times 2=739769049169920 \approx 2^{49.4}$,
  \item $7 \times 100 \times 147 \times 81 \times 24 \times 2 \times 3 \times 80 \times 84 \times 28 \times 4=903337795584000 \approx 2^{49.7}$,
  \item $7 \times 100 \times 147 \times 87 \times 35 \times 1 \times 3 \times 80 \times 78 \times 27 \times 1=158369767920000 \approx 2^{47.2}$.
\end{itemize}

Note that the largest number of the above combinations is smaller than $2^{50}$.\\
We remark that, in Proposition 3, the cardinalities of the candidate values for $x_{13}, x_{12}, x_{7}, x_{6}$ are $7,2^{7}-2^{4}, 3$ and $2^{7}-2^{5}$, respectively; but in the experiments with respect to Conjecture 3 , the cardinalities of the candidate values for $x_{12}$ and $x_{6}$ are less than $2^{7}-2^{4}$ and $2^{7}-2^{5}$. This is because, when running Algorithm 1 of [14], we traverse $x_{13}$ first, then $x_{12}, x_{11}, \cdots, x_{2}, x_{1}$. So the value of $x_{i}$ will have an impact on the value of $x_{j}(i>j)$.

Note that the complexity of the exhaustive method in dimension 8 is $2^{8 \times 13}=2^{104}$, while the complexity of the QAM method is $2^{50}$, which is why we consider the QAM method to be rather valuable even now.

\section*{5 A new model}
Algorithm 1 in [14] handles the last two columns of $H_{8}\left(x_{13}, x_{12}, \ldots, x_{2}\right.$ and $\left.x_{1}\right)$ one by one. In this section, we propose a new model which can handle the last two columns together. Our model can aviod some redundant computation. However, up to now, no new APN functions have been found by using this model, as a matter of fact, we haven't realized the new model. We want to study this model further in the future.

We need the following results to prove Theorem 6.

Definition 3 Let $H \in \mathbb{F}_{2^{n}}^{m \times k}(m, k \leq n)$. $H$ is called proper if every nonzero linear combination of the $m$ rows of $H$ has rank (over $\mathbb{F}_{2}$ ) at least $k-1$.

Theorem 4 (Theorem 4 of 14) Let $A=\left(a_{i, j}\right) \in \mathbb{F}_{2^{n}}^{r \times c}(1 \leq r<c \leq n)$ such that $a_{i, j}=a_{j, i}$ and $a_{i, i}=0$ for $1 \leq i, j \leq r$. Let $A[\cdot, k]$ be the $k$-th column of $A$, and $b=\sum_{k=1}^{c} \lambda_{k} A[\cdot, k]$, where $\lambda_{k} \in \mathbb{F}_{2}$ for $1 \leq k \leq c$. Assume $t=\operatorname{Rank}_{\mathbb{F}_{2}}\{b[1], b[2], \ldots, b[r]\}$ and $\left(\lambda_{1}, \ldots, \lambda_{c}\right) \neq 0$. If $A$ is proper, then we have:\\
i) If $\left(\lambda_{r+1}, \ldots, \lambda_{c}\right)=0$, then $t=r-1$;\\
ii) If $\left(\lambda_{r+1}, \ldots, \lambda_{c}\right) \neq 0$, then $t=r$.

Based on the definitions of QAM and proper, it is easy to prove the following lemma.

Lemma 5 Suppose $H$ is an $n \times n$ matrix over $\mathbb{F}_{2^{n}}$, and $A$ is a $m \times k$ submatrix of $H$. If $H$ is a QAM, then A must be proper.

Proof Suppose $A$ is not proper, then there must exist a nonzero linear combination of the $m$ rows of $A$ that has rank less than $k-1$, which implies that there exist a nonzero linear combination of the $n$ rows of $H$ that has rank less than $n-1$. In other words, $A$ is not proper implies that $H$ is not a QAM. So $A$ must be proper.

Now we are ready to prove the following theorem.

Theorem 6 Let $H$ be the corresponding $Q A M$ of a quadratic APN function on $\mathbb{F}_{2^{n}}$. Define an $n \times n$ symmetric matrix $H_{n}$ such that submatrix $\left(H_{n}, 1,1, n-2, n-2\right)=\operatorname{submatrix}(H, 1,1, n -2, n-2)$ (The left upper $(n-2) \times(n-2)$ submatrix of $H$ ) and $H_{n}[i, i]=0$ for $n-1 \leq i \leq n$. Let $H_{n}[i, 1 \rightarrow n-2]=\left(H_{n}[i, 1], H_{n}[i, 2], \cdots, H_{n}[i, n-2]\right)$. Then $H_{n}$ is a QAM if and only if


\begin{equation*}
\sum_{i=1}^{n-1} \lambda_{i}\left(\mu_{1} H_{n}[i, n-1]+\mu_{2} H_{n}[i, n]\right) \in \mathbb{F}_{2^{n}} \backslash \operatorname{Span}\left(\sum_{i=1}^{n-1} \lambda_{i} H_{n}[i, 1 \rightarrow n-2]\right) \tag{2}
\end{equation*}


for all $\left(\lambda_{1}, \lambda_{2}, \cdots, \lambda_{n-1}, \mu_{1}, \mu_{2}\right) \neq(0,0, \cdots, 0,0,1,1,0), \quad\left(\lambda_{1}, \lambda_{2}, \cdots, \lambda_{n-1}\right) \in \mathbb{F}_{2}^{n-1} \backslash\{0\} \quad$ and $\left(\mu_{1}, \mu_{2}\right) \in \mathbb{F}_{2}^{2} \backslash\{0\}$.

Proof If $H_{n}$ is a QAM, then (2) is obvious true. Next we suppose that (2) is true, and we prove that $H_{n}$ is a QAM.

\begin{enumerate}
  \item When $\left(\lambda_{1}, \lambda_{2}, \cdots, \lambda_{n-1}, \mu_{1}, \mu_{2}\right)=(0,0, \cdots, 0,0,1,1,0)$, (2) becomes $H_{n}[n-1, n-1] \in \mathbb{F}_{2^{n}} \backslash \operatorname{Span}\left(H_{n}[n-1,1 \rightarrow n-2]\right)$, which is impossible since $H_{n}[n -1, n-1]=0$ is a known value.
  \item According to the definition of $H_{n}$, it is symmetric and the elements in its main diagonal are all zeros.
  \item Equation (2) implies that every nonzero linear combination of the $n$ rows of $H_{n}$ has rank $n-1$ over $\mathbb{F}_{2}$.
\end{enumerate}

According to 1 ), 2) and 3), $H_{n}$ is a QAM.\\
Part 3) is not obvious, we will prove it in the following.\\
Let $A=\operatorname{submatrix}\left(H_{n}, 1,1, n-2, n-2\right)$, then according to Lemma 5, the ( $\left.n-2\right) \times(n-2)$ matrix $A$ is proper and


\begin{equation*}
\operatorname{Rank}_{\mathbb{F}_{2}}\left(\sum_{i=1}^{n-2} \lambda_{i} A[i]\right) \geq n-3 \text { for all }\left(\lambda_{1}, \lambda_{2}, \cdots, \lambda_{n-2}\right) \in \mathbb{F}_{2}^{n-2} \backslash\{0\} . \tag{3}
\end{equation*}


Let $\left(\mu_{1}, \mu_{2}\right)=(1,0)$ and $B=\operatorname{submatrix}\left(\underset{n-2}{H_{n-2}}, 1,1, n-2, n-1\right)$. Then according to (2) and (3), we have $\sum_{i=1} \lambda_{i} H_{n}[i, n-1] \in \mathbb{F}_{2^{n}} \backslash \operatorname{Span}\left(\sum_{i=1} \lambda_{i} A[i]\right)$ for all $\left(\lambda_{1}, \lambda_{2}, \cdots, \lambda_{n-2}\right) \in \mathbb{F}_{2}^{n-2} \backslash\{0\}$, which implies that $\operatorname{Span}\left(\sum_{i=1}^{n-2} \lambda_{i} B[i]\right) \geq n-2$ for all $\left(\lambda_{1}, \lambda_{2}, \cdots, \lambda_{n-2}\right) \in \mathbb{F}_{2}^{n-2} \backslash\{0\}$. So $B$ is an $(n-2) \times(n-1)$ proper matrix. Let $C=\operatorname{submatrix}\left(H_{n}, 1,1, n-1, n-1\right)$, then according to Theorem 4 and the symmetric property of $H_{n}$ we can conclude that the $(n-1) \times(n-1)$ matrix $C$ is proper and


\begin{equation*}
\operatorname{Rank}_{\mathbb{F}_{2}}\left(\sum_{i=1}^{n-1} \lambda_{i} C[i]\right)=n-2 \tag{4}
\end{equation*}


for all $\left(\lambda_{1}, \lambda_{2}, \cdots, \lambda_{n-1}\right) \in \mathbb{F}_{2}^{n-1} \backslash\{0\}$.\\
Let $D=\operatorname{submatrix}\left(H_{n}, 1,1, n-1, n\right)$, and choosing $\left(\mu_{1}, \mu_{2}\right)=(1,1)$ and $\left(\mu_{1}, \mu_{2}\right)=(0,1)$ respectively, then according to (2) and (4), we can conclude that $D$ is a proper matrix. Similar as the above proof, we can conclude that

$$
\operatorname{Rank}_{\mathbb{F}_{2}}\left(\sum_{i=1}^{n} \lambda_{i} H_{n}[i]\right)=n-1
$$

for all $\left(\lambda_{1}, \lambda_{2}, \cdots, \lambda_{n}\right) \in \mathbb{F}_{2}^{n} \backslash\{0\}\left(H_{n}[i]\right.$ denotes the $i$-th row of $H_{n}$, and similar as $A[i], B[i]$ and $C[i]$ ).

Algorithm 1 in [14] handles the last two columns one by one, while Theorem 6 provides a theoretical foundation to handle them together. One advantage to handle the last two columns together is that it can avoid some redundant computation. For example, in our model, when traversing the last two columns of $H_{8}, x_{13}$ and $x_{7}$ must satisfy

\[
\left\{\begin{array}{l}
x_{13} \in \mathbb{F}_{2^{8}} \backslash \operatorname{Span}\left(g^{34}, g^{81}, g^{83}, g^{170}, g^{106}\right)  \tag{5}\\
x_{7} \in \mathbb{F}_{2^{8}} \backslash \operatorname{Span}\left(g^{34}, g^{81}, g^{83}, g^{170}, g^{106}\right) \\
x_{13}+x_{7} \in \mathbb{F}_{2^{8}} \backslash \operatorname{Span}\left(g^{34}, g^{81}, g^{83}, g^{170}, g^{106}\right)
\end{array}\right.
\]

While in Algorithm 1 in [14], $x_{13}$ and $x_{7}$ must satisfy

\[
\left\{\begin{array}{l}
x_{13} \in \mathbb{F}_{2^{8}} \backslash \operatorname{Span}\left(g^{34}, g^{81}, g^{83}, g^{170}, g^{106}\right)  \tag{6}\\
x_{7} \in \mathbb{F}_{2^{8}} \backslash \operatorname{Span}\left(g^{34}, g^{81}, g^{83}, g^{170}, g^{106}, x_{13}\right)
\end{array}\right.
\]

Note that in (5), we need to compute $\mathbb{F}_{2^{8}} \backslash \operatorname{Span}\left(g^{34}, g^{81}, g^{83}, g^{170}, g^{106}\right)$ only once, while in (6), we need to compute $\mathbb{F}_{2^{8}} \backslash \operatorname{Span}\left(g^{34}, g^{81}, g^{83}, g^{170}, g^{106}, x_{13}\right)$ for many candidates of $x_{13}$. Similarly property holds for all the values of the last two columns. It is easy to see that the new model can avoid some redundant computation.

Acknowledgements Yuyin Yu is supported by the National Key R\&D Program of China (Grant No. 2021YFB3100200) and the GuangDong Basic and Applied Basic Research Foundation (Grant No. 2021A1515011904). We thank the anonymous reviewers for their comments.

\section*{References}
\begin{enumerate}
  \item Beierle, C., Brinkmann, M., Leander, G.: Linearly self-equivalent APN permutations in small dimension. IEEE Trans. Inf. Theory 67(7), 4863-4875 (2021)
  \item Budaghyan, L.: Construction and Analysis of Cryptographic Functions, pp. 1-168. Springer (2014) ISBN 978-3-319-12990-7
  \item Beierle, C., Leander, G.: New instances of quadratic APN functions. IEEE Trans. Inf. Theory 68(1), 670-678 (2022)
  \item Bosma, W., Cannon, J., Playoust, C.: The Magma algebra system I: The user language[J]. J. Symb. Comput. 24(3-4), 235-265 (1997)
  \item Browning, K., Dillon, J.F., McQuistan, M.T., Wolfe, A.J.: An APN permutation in dimension six. Contemp. Math. 58, 33-42 (2010)
  \item Canteaut, A., Couvreur, A., Perrin, L.: Recovering or testing extended-affine equivalence. IEEE Trans. Inf. Theory. \href{https://doi.org/10.1109/TIT.2022.3166692}{https://doi.org/10.1109/TIT.2022.3166692}
  \item Carlet, C.: Boolean Functions for Cryptography and Coding Theory, pp. 369-424. Cambridge University Press (2021) ISBN 978-1-108-47380-4 (Chapter 11)
  \item Carlet, C., Charpin, P., Zinoviev, V.A.: Codes, bent functions and permutations suitable for des-like cryptosystems. Des Codes Cryptogr. 15(2), 125-156 (1998)
  \item Edel, Y., Pott, A.: A new almost perfect nonlinear function which is not quadratic. Adv. Math. Commun. 3(1), 59-81 (2009)
  \item Kaleyski, N.: Deciding EA-equivalence via invariants. Cryptograp. Commun. 27, 1-20 (2021)
  \item Taniguchi, H.: On some quadratic APN functions. Des. Codes Crypt. 87, 1973-1983 (2019)
  \item Weng, G., Tan, Y., Gong, G.: On quadratic almost perfect nonlinear functions and their related algebraic object. In: Workshop on Coding and Cryptography, WCC (2013)
  \item Yoshiara, S.: Equivalences of quadratic APN functions. J. Algebr. Comb. 35, 461-475 (2011)
  \item Yu, Y., Wang, M., Li, Y.: A matrix approach for constructing quadratic APN functions. Des. Codes Cryptog. 73, 587-600 (2014)
\end{enumerate}

Publisher's note Springer Nature remains neutral with regard to jurisdictional claims in published maps and institutional affiliations.


\end{document}